% !TITLE 鹊踏枝·谁道闲情抛掷久
% !AUTHOR 冯延巳
% !DYNASTY 五代十国
% !GENRE 词
% !ORDER 7

\begin{poem}
谁道闲情\textsuperscript{1}抛掷\textsuperscript{2}久？每到春来，惆怅\textsuperscript{3}还依旧。
日日花前常病酒\textsuperscript{4}，不辞\textsuperscript{5}镜里朱颜\textsuperscript{6}瘦。

河畔青芜\textsuperscript{7}堤上柳，为问\textsuperscript{8}新愁\textsuperscript{9}，何事年年有？
独立小桥风满袖，平林\textsuperscript{10}新月人归后。
\end{poem}

\begin{annotations}
\notetext{1}{闲情：无端的愁绪、说不清道不明的烦恼。不是具体的愁，而是一种弥漫在春天里的莫名忧伤。}
\notetext{2}{抛掷：扔掉、摆脱。谁说我已经把闲愁丢掉了？一个反问，说明根本没有丢掉。}
\notetext{3}{惆怅：失意、伤感。春天一来，忧伤又恢复了原来的样子，一点没减少。}
\notetext{4}{病酒：因饮酒过量而身体不适。天天在花前喝酒喝到伤身——宁可伤害身体，也要排遣愁绪。}
\notetext{5}{不辞：不推辞、不躲避。不介意镜子里的自己容颜消瘦。为了解愁，不惜代价。}
\notetext{6}{朱颜：红润的面容，指青春美貌。镜中看到自己一天天憔悴，却说“不辞”，这是以倔强写脆弱。}
\notetext{7}{青芜：青色的野草，丛生的荒草。芜（wú），杂草。河畔蔓草、堤上垂柳，都是春天新生的愁绪的化身。}
\notetext{8}{为问：试问、请问。“为问新愁”——我想问一问这新来的愁绪。向自然发问，实则问自己。}
\notetext{9}{新愁：新的忧愁。年年都有新愁，和旧愁混在一起，永远排遣不尽。}
\notetext{10}{平林：平地上的树林。月上林梢，人们都已归家，只有词人独自站在小桥上，袖子灌满了风——“独立”和“人归后”构成强烈反差。}
\end{annotations}

\begin{appreciation}
五代词的收尾，我们读冯延巳。他是南唐的宰相，李煜父亲的重臣，排在本书最后，因为他的词是五代词的句号，也是两宋词的起笔。北宋的晏殊、欧阳修都从他这里学东西。顺便记一个知识点：鹊踏枝这个词牌，就是欧阳修写“庭院深深深几许”用的蝶恋花——同一个曲调，两个名字。

这首词写的是“闲情”。闲情不是失恋，不是失业，是无缘无故的愁，用现在的话说，就是没来由的emo。“谁道闲情抛掷久？”谁说我把闲情扔掉了很久？一个反问，答案马上接上来：“每到春来，惆怅还依旧。”春天一到，惆怅又回来了，跟去年一模一样，根本扔不掉。

“日日花前常病酒，不辞镜里朱颜瘦。”天天在花前喝酒，喝到身体不舒服——病酒，就是喝伤了的难受。不辞，是不在乎：镜子里看到自己憔悴了，不在乎。这种豁出去的样子，是用倔强写脆弱，跟李煜“罗衾不耐五更寒”的硬撑，是同一种撑法。“河畔青芜堤上柳，为问新愁，何事年年有？”河边的青草、堤上的柳树，年年春天都发新绿，他问：为什么新愁也跟野草一样，年年春风吹又生？

最后两句是词史上最孤独的画面之一：“独立小桥风满袖，平林新月人归后。”一个人站在小桥上，风灌满了袖子；树林尽头新月升起，所有人都回家了，他还站着。没有一句愁字，愁却站成了一个人形。

李煜的愁有名字，叫亡国；冯延巳的愁没有名字。他亲眼看着南唐一天天衰败下去，却什么也说不出口——有些愁，是说不清来由的，说出来就变味了。诗里的愁，向来有事可愁；词从冯延巳开始，才第一次写“没有事的愁”，这一写，就写进了宋词这条大河里。

你以后要是哪天无缘无故不高兴，不用找原因，也不用逼自己开心。春天会来，惆怅也会来，都是正常的事。饿了就吃，困了就睡，过两天自己就好了。
\end{appreciation}
